% GDZ PPN236018647,native507;HKUST508;original diary9v,service20.
% Complete manuscript entries corresponding to printed114-120.
\clearpage\SourceLeaf{507}\NoSourceHead
\begingroup\fontsize{10}{15}\selectfont\itshape
\newcommand{\DiaryHashSeven}{\tikz[baseline=-.3ex,x=.7em,y=.7em]{\draw[line width=.4pt](.08,-.65)--(.48,.8);\draw[line width=.4pt](.48,-.7)--(.78,.85);\draw[line width=.45pt](-.05,.35)..controls(.4,.4)and(.7,.25)..(1.05,.24);\draw[line width=.45pt](-.08,.02)--(.98,-.11);}}
\noindent\begin{minipage}[t]{.15\linewidth}\vspace{0pt}
\tikz[x=.11em,y=.11em]{
\draw[line width=.4pt] (7,40)..controls(14,50)and(25,55)..(34,55)..controls(48,54)and(56,47)..(56,36)..controls(56,27)and(49,17)..(43,10)..controls(36,3)and(31,0)..(25,0)..controls(12,0)and(6,6)..(3,16)..controls(0,25)and(0,31)..(4,37);
\draw[line width=.4pt] (17,39)..controls(21,37)and(23,32)..(28,28)..controls(33,25)and(37,25)..(42,24);
\draw[line width=.4pt] (18,38)..controls(21,30)and(25,20)..(28,9);
\draw[line width=.45pt] (30,43)..controls(29,34)and(26,26)..(23,19)..controls(21,15)and(19,14)..(18,13);
\draw[line width=.4pt] (32,43)..controls(33,31)and(35,23)..(39,13);
\draw[line width=.4pt] (38,36)--(14,28);
\draw[line width=.4pt] (38,36)..controls(34,28)and(31,18)..(29,9);
\draw[line width=.4pt] (42,23)--(18,13);
\draw[line width=.65pt] (14,25)..controls(22,20)and(30,17)..(38,12);
\draw[line width=.5pt] (16,39)..controls(12,39)and(13,44)..(16,45)..controls(20,46)and(21,40)..(16,39);
\draw[line width=.45pt] (30,43)..controls(26,44)and(29,50)..(32,50)..controls(36,50)and(36,44)..(32,43)--(30,43);
\draw[line width=.45pt] (38,36)..controls(34,37)and(38,42)..(41,42)..controls(45,40)and(42,36)..(38,36);
\draw[line width=.5pt] (42,24)..controls(41,28)and(43,32)..(45,31)..controls(49,29)and(48,23)..(45,22)--(42,23)--(42,24);
\draw[line width=.5pt] (38,12)..controls(36,13)and(38,16)..(41,15)..controls(46,12)and(42,8)..(39,10)--(38,12);
\draw[line width=.55pt] (28,9)..controls(24,9)and(24,3)..(27,2)..controls(32,2)and(33,8)..(29,9)--(28,9);
\draw[line width=.5pt] (18,13)..controls(14,17)and(11,11)..(12,8)..controls(15,5)and(20,10)..(19,13)--(18,13);
\draw[line width=.6pt] (14,28)..controls(10,33)and(5,28)..(5,23)..controls(7,18)and(14,22)..(14,25)--(14,28);
}
\end{minipage}\hfill\begin{minipage}[t]{.82\linewidth}\vspace{0pt}
\noindent Nov. 30, Felix fuit dies quo multitudinem\\
classium formar binar. per triplicem methodum\\
assignare largitum est nobis puta 1) per prod. infin.
\end{minipage}
\par\noindent 2) per aggregatum infinitum, 3) per aggregatum fini-\\
tum cotangentium seu logarithm. sinuum. Brun.

\medskip
\noindent\DiaryHashSeven\quad Dec. 3. Methodum quartam \raisebox{1.15ex}{\small ex} omnibus simplicissimam\\
deteximus pro dett. negativis ex sola multit. nume-\\
ror \(\rho;\rho'\) etc. petitam si \(Ax+\rho,Ax+\rho'\) etc. sunt\\
formae lineares divisor for \(\Box+D\) \raisebox{.45ex}{.}\hfill Ibid.
\par\smallskip\noindent\rule{\linewidth}{.3pt}

\noindent\tikz[baseline=-.25ex,x=.6em,y=.6em]{
\draw[line width=.45pt](.1,-.5)--(.75,.67);\draw[line width=.35pt](.28,-.55)--(.85,.55);\draw[line width=.45pt](.12,.55)--(.83,-.2);\draw[line width=.4pt](0,.25)--(.72,-.42);
}\quad Impossibile esse ut sectio circuli ad aequationes inferi-\\
ores quam theoria nostra suggerit reducatur demon-\\
stratum\hfill Brunov. Apr. 6
\par\smallskip\noindent\rule{\linewidth}{.3pt}

\noindent\underline{Iisdem diebus Pascha Iudaeorum per methodum novam}\\
\underline{determinare docuimus}\quad (Apr. 1)
\par\smallskip\noindent\rule{\linewidth}{.3pt}

\noindent\tikz[baseline=-.25ex,x=.65em,y=.65em]{\draw[line width=.5pt](0,.13)--(.75,.2);\draw[line width=.45pt](.2,-.25)--(.48,.63);\draw[line width=.4pt](.15,.48)--(.64,-.12);}\quad\underline{Methodus quinta theorema fundamentale demonstrandi}\\
\underline{se obtulit adiumento theorematis elegantissimi theoriae sectionis circuli}
\par\noindent\hspace*{12em}puta
\[
\begin{array}{r@{\;}cccc@{\;}l|l}
\displaystyle\sum\left.\begin{matrix}\mathit{sin}\\\mathit{cos}\end{matrix}\right\}\frac{nn}{a}P=
&\multicolumn{1}{c|}{\begin{matrix}+\surd a\\+\surd a\end{matrix}}
&\multicolumn{1}{c|}{\begin{matrix}0\\+\surd a\end{matrix}}
&\multicolumn{1}{c|}{\begin{matrix}0\\0\end{matrix}}
&\begin{matrix}+\surd a\\0\end{matrix}&
&\begin{matrix}\text{\itshape substituendo pro }n\\\text{\itshape omnes numeros}\end{matrix}\\[5pt]
\text{\itshape prout }a\equiv&0&1&2&3&(\mathrm{mod.}\ 4)&\text{\itshape a }0\text{\itshape usque ad }a-1
\end{array}
\]
\noindent\hfill Brunsv. Mai. medio
\par\smallskip\noindent\rule{\linewidth}{.3pt}

\noindent\DiaryHashSeven\quad\underline{Methodus nova simplicissima expeditissi}\\
\underline{ma \tikz[baseline=(t.base)]{\node[inner sep=0pt,outer sep=0pt] (t) {\strut moti};\draw[line width=.65pt] ([yshift=.1ex]t.west)..controls([yshift=-.1ex]t.center)..([xshift=.08em,yshift=.03ex]t.east);\draw[line width=.4pt] ([xshift=.1em,yshift=-.4ex]t.west)--([xshift=.04em,yshift=-.35ex]t.east);} elementa orbitarum corporum coele}\\
\underline{stium investigandi}.\hfill Brunsv. Sept. m

\medskip
\noindent\DiaryHashSeven\quad Theoriam motus Lunae aggressi sumus\hfill Aug.
\endgroup
